\documentclass{article}
% Fixture for the claim-forward module (latex-paper-en).
% Cases A-E are positives (one per CF-* code); Cases F-H are boundaries that
% must stay silent. Written for evals/evals.json id 24 and
% tests/skills/latex_paper_en/test_claim_forward.py.
\begin{document}

\begin{abstract}
We present a lightweight tracker that achieves 12\% higher accuracy than the strongest baseline on three indoor benchmarks.
The tracker runs at 210 frames per second on a single GPU.
\end{abstract}

\section{Introduction}
% Case A: CF-DISCLAIM — paragraph opens by denying scope before stating the claim.
We do not claim to solve general scene understanding.
Instead, this paper presents a lightweight tracker that achieves 12\% higher accuracy on indoor benchmarks.

\section{Contributions}
% Case C: CF-CAVEAT-POS — limitation sentence ahead of the claim it qualifies.
Although the evaluation is restricted to indoor scenes, the conclusions should be read with that scope in mind.
We introduce a tracker that reduces latency by 40\% relative to the best published method.

\section{Method}
% Case D: CF-HEDGE-STACK — four hedges on one claim sentence.
Our approach may potentially improve robustness to some extent in some cases.
The matching module uses a fixed window of eight frames.

\section{Results}
% Case B: CF-SELFWEAK — self-weakening collocations on the authors' own result.
Our method achieves 94.1\% accuracy on the standard split.
Regrettably, our method still lags far behind the oracle on the long-tail split.

% Case H: boundary — bare "only" and "limited" are not self-weakening.
Our method achieves 3\% higher accuracy than the best baseline.
It only requires a single GPU and has a limited memory footprint of 2~GB.

\section{Discussion}
% Case F: boundary — cited prior work may be described as falling short.
Prior work~\cite{tracker2019} falls short of real-time rates on commodity hardware.
These methods also suffer from serious drift over long sequences.

\section*{Limitations}
% Case G: boundary — a Limitations section is where limitations belong.
Our evaluation does not cover outdoor scenes with strong illumination change.
We show results on three indoor benchmarks only.

\section{Conclusion}
We presented a lightweight tracker and showed a 12\% accuracy gain on three indoor benchmarks.
The design keeps the per-frame cost below five milliseconds.

% Case E: CF-CLOSE-NEG — final paragraph ends on a negative judgment with no direction.
Our framework achieves real-time rates on a single GPU.
However, it does not handle streaming inputs with variable frame rates.

\end{document}
